% BANKS-SOURCE: pdf=253 print=243 kind=solution-section id=chapter11-implicit-solutions
\begin{bankssolutions}{Solutions to Implicit Exercises for Chapter 11}

% BANKS-SOURCE: pdf=253 print=243 kind=solution id=I-CH11-001
\BanksImplicitSolution{I-CH11-001}
The background metric is fixed throughout; the functional integral runs only
over the matter field.  We use the mostly-plus signature adopted here,
$\eta_{\mu\nu}=\operatorname{diag}(-,+,\ldots,+)$, continue along
$t=-\ii\tau$, and assume that the couplings and the background admit the same
analytic continuation as the metric.  The static Hamiltonian is bounded below
and has a normalizable ground state.  Throughout, $D$ denotes the spacetime
dimension and remains distinct from any regulator parameter.

On the Euclidean slice, $\dd z=-\ii\dd\tau$, so the holomorphic metric first
restricts to the real positive-definite line element
\[
 \dd s_{\mathbb C}^2|_{M_E}
 =N^2\dd\tau^2+h_{ij}\dd x^i\dd x^j.
\]
With $g_E=g_{\mathbb C}|_{M_E}$, the positive-definite Euclidean line element
is
\[
 \dd s_E^2
 =N^2\dd\tau^2+h_{ij}\dd x^i\dd x^j.
\]
The assumptions $N>0$ and $h_{ij}>0$ make this metric positive definite.  The
coordinate volume factors on the two real slices are
\[
 \sqrt{-g_L}=N\sqrt h,
 \qquad
 \sqrt{g_E}=N\sqrt h.
\]
Since $t=-\ii\tau$, one has $\dd t=-\ii\dd\tau$ and
$\partial_t=\ii\partial_\tau$.  Direct substitution in the Lorentzian action
gives
\begin{align*}
 S_L\big|_{t=-\ii\tau}
 &=(-\ii)\int\dd\tau\,\dd^{D-1}\mathbf x\,N\sqrt h\\
 &\qquad {}\times\left[
 -\frac{(\partial_\tau\phi)^2}{2N^2}
 -\frac12h^{ij}\partial_i\phi\partial_j\phi
 -\frac12m^2\phi^2-V_{\mathrm{int}}(\phi)
 \right]\\
 &=\ii S_E,
\end{align*}
where
\[
 S_E=\int_{M_E}\dd\tau\,\dd^{D-1}\mathbf x\,N\sqrt h\left[
 \frac{(\partial_\tau\phi)^2}{2N^2}
 +\frac12h^{ij}\partial_i\phi\partial_j\phi
 +\frac12m^2\phi^2+V_{\mathrm{int}}(\phi)
 \right].
\]
It follows at once that $\ee^{\ii S_L}$ continues to $\ee^{-S_E}$.

With a regulator understood, define
\[
 \mathcal Z_E[J]=\int_{\mathcal B}\mathcal D\phi\,
 \exp\left[-S_E[\phi]
 +\ii\int_{M_E}\dd^D x\,\sqrt{g_E}\,J(x)\phi(x)\right],
 \qquad
 Z_E[J]=\frac{\mathcal Z_E[J]}{\mathcal Z_E[0]}.
\]
The boundary conditions $\mathcal B$ specify the state.  The normalized
Schwinger functions are
\[
 G_E^{(n)}(x_1,\ldots,x_n)
 =\left.\frac{1}{\ii^n\sqrt{g_E(x_1)}\cdots\sqrt{g_E(x_n)}}
 \frac{\delta^n Z_E[J]}
 {\delta J(x_1)\cdots\delta J(x_n)}\right|_{J=0}.
\]
Euclidean evolution over semi-infinite time intervals, with regular
finite-action boundary data, projects onto the lowest-energy state associated
with the static Killing field.  A smooth Euclidean circle
$\tau\sim\tau+\beta$ gives
$\mathcal Z_E[0]=\Tr\ee^{-\beta H}$ and prepares the thermal density matrix
$\rho_\beta=\ee^{-\beta H}/\Tr\ee^{-\beta H}$.  Its bosonic correlators are
periodic in $\tau$ and continue to KMS correlators.

To derive the continuation, write $K=H-E_0$ and choose
$K\ket{\Omega}=0$.  Write $K\ket{r}=E_r\ket{r}$, with $E_r\geq0$.
Euclidean evolution is
\[
 \phi_E(\tau,\mathbf x)
 =\ee^{K\tau}\phi(0,\mathbf x)\ee^{-K\tau}.
\]
For $\tau_1>\tau_2>\cdots>\tau_n$, Euclidean ordering therefore gives
\begin{align*}
 G_E^{(n)}(1,\ldots,n)
 ={}&\bra{\Omega}\phi_1
 \ee^{-K(\tau_1-\tau_2)}\phi_2
 \ee^{-K(\tau_2-\tau_3)}\cdots
 \ee^{-K(\tau_{n-1}-\tau_n)}\phi_n\ket{\Omega},
\end{align*}
where $\phi_a=\phi(0,\mathbf x_a)$.  Inserting complete sets of energy
eigenstates turns this expression into sums of terms containing
\[
 \ee^{-E_{r_1}(\tau_1-\tau_2)}
 \ee^{-E_{r_2}(\tau_2-\tau_3)}\cdots
 \ee^{-E_{r_{n-1}}(\tau_{n-1}-\tau_n)}.
\]
Every such term is holomorphic when
$\operatorname{Re}(\tau_a-\tau_{a+1})>0$.  The same statement follows from a
continuous spectral measure when the spectrum has a continuous part.

Write $\zeta$ for complexified Euclidean time.  Let $\pi$ order the Lorentzian
times,
$t_{\pi(1)}>t_{\pi(2)}>\cdots>t_{\pi(n)}$.  For $\epsilon>0$, set
\[
 \zeta_{\pi(a)}=\ii t_{\pi(a)}+(n-a)\epsilon.
\]
These points remain in the ordered holomorphy domain because
$\operatorname{Re}(\zeta_{\pi(a)}-\zeta_{\pi(a+1)})=\epsilon$.  Continuing the
spectral factors yields
\[
 \ee^{-E_r(\zeta_{\pi(a)}-\zeta_{\pi(a+1)})}
 =\ee^{-\epsilon E_r}
  \ee^{-\ii E_r(t_{\pi(a)}-t_{\pi(a+1)})}.
\]
Equivalently, the associated Lorentzian times approach the real axis as
$t_{\pi(a)}\mapsto t_{\pi(a)}-\ii(n-a)\epsilon$.  This is the ordered
$\ii\epsilon$ prescription.
Taking the boundary value gives
\[
 \boxed{
 \bra{\Omega}T\{\phi_L(t_1,\mathbf x_1)\cdots
 \phi_L(t_n,\mathbf x_n)\}\ket{\Omega}
 =\lim_{\epsilon\downarrow0}
 G_E^{(n)}(\zeta_{\pi(1)},\mathbf x_{\pi(1)};\ldots;
 \zeta_{\pi(n)},\mathbf x_{\pi(n)}).}
\]
The positive real gaps damp high energies and prescribe the side from which
the singular Lorentzian boundary is reached.  Coincident-point expressions
are understood as distributions after renormalization.

The static metric used above supplies a concrete sufficient geometry.  More
intrinsically, $M_L$ and $M_E$ are fixed sets of antiholomorphic real
structures on the same $M_{\mathbb C}$, and the holomorphic metric restricts
to a real Lorentzian metric on $M_L$ and a real definite metric on $M_E$; with
the mostly-plus convention, the positive Euclidean metric is
$g_E=g_{\mathbb C}|_{M_E}$.  The slices must be connected through a domain
where the complex metric stays nondegenerate and the continued Green functions
remain clear of singularities.  A globally hyperbolic Lorentzian region gives
a well-posed field evolution.  If the Euclidean section closes at a Lorentzian
Killing horizon, smoothness fixes the period of $\tau$ by removing the conical
singularity.

For the field theory, the metric, couplings, and insertions need enough
analyticity to extend into that common domain.  The Euclidean kinetic operator
is elliptic, and its boundary conditions give a unique Green function after
zero modes are treated.  Convergence or a regulator defines the functional
integral.  Reflection positivity under Euclidean time reflection supplies a
positive Hilbert-space inner product after continuation, while locality and
regularity support the reconstruction.  These conditions explain the
content of Banks's phrase ``the appropriate Euclidean manifold'': its geometry
and regularity select the Lorentzian state whose correlators appear as boundary
values.

For the flat-space check, put
$N=1$, $h_{ij}=\delta_{ij}$, and
$\omega_{\mathbf p}=\sqrt{\mathbf p^2+m^2}$.  The free Euclidean propagator is
\begin{align*}
 G_E(\tau,\mathbf x)
 &=\int\frac{\dd p_E^0\,\dd^{D-1}\mathbf p}{(2\pi)^D}
 \frac{\ee^{\ii p_E^0\tau+\ii\mathbf p\cdot\mathbf x}}
 {(p_E^0)^2+\omega_{\mathbf p}^2}\\
 &=\int\frac{\dd^{D-1}\mathbf p}{(2\pi)^{D-1}}
 \frac{\ee^{\ii\mathbf p\cdot\mathbf x}}
 {2\omega_{\mathbf p}}\ee^{-\omega_{\mathbf p}|\tau|}.
\end{align*}
For $t>0$, continue the positive-$\tau$ branch through
$\tau=\epsilon+\ii t$.  For $t<0$, continue the negative-$\tau$ branch through
$\tau=-\epsilon+\ii t$.  Both boundary values combine into
\[
 G_F(t,\mathbf x)
 =\int\frac{\dd^{D-1}\mathbf p}{(2\pi)^{D-1}}
 \frac{\ee^{\ii\mathbf p\cdot\mathbf x}}
 {2\omega_{\mathbf p}}\ee^{-\ii\omega_{\mathbf p}|t|}.
\]
Here $p\mathbin{\cdot}x=-p^0t+\mathbf p\mathbin{\cdot}\mathbf x$ and
$p^2=-(p^0)^2+\mathbf p^2$, so the positive-frequency phase is
$\ee^{\ii p\mathbin{\cdot}x}$ and the mass shell is $p^2=-m^2$.
Equivalently,
\[
 G_F(t,\mathbf x)
 =\int\frac{\dd p^0\,\dd^{D-1}\mathbf p}{(2\pi)^D}
 \frac{-\ii\,\ee^{\ii p\mathbin{\cdot}x}}
 {p^2+m^2-\ii0}.
\]
For positive $t$ the pole prescription selects
$\ee^{-\ii\omega_{\mathbf p}t}$; for negative $t$ it selects
$\ee^{\ii\omega_{\mathbf p}t}$.  This is the positive-frequency Feynman
two-point function and checks the continuation, its time ordering, and its
$\ii0$ prescription.

\end{bankssolutions}
